Benchmark scores for network intrusion detectors hide problems that appear in deployment. We audit a decision tree on 54 NFStream flow features from labels to live blocking on an emulated Open vSwitch gateway, and condense the audit into four checks. \emph{Labels:} background, i.e.\ traffic outside the attack, makes up 90.9--98.5\% of the flows in six of the nine CIC-IoT2023 attack captures. Trained on file-name labels, the tree flags 71--88\% of the background flows in unseen attack captures, against 0.5--5.4\% with evidence labels; on IoT-23 it flags 23--47\% of unseen benign devices' flows, against 0--17.5\% with the authors' Zeek labels. \emph{Splits:} on CIC-IoT2023 a random split hides that brute-force recall falls from 1.00 to 0.05 when the attack moves from SSH passwords to RTSP stream paths; this case (three host pairs) is most of the random--temporal gap. On IoT-23 the gap comes from classes missing from one split half. \emph{Live:} matching live to offline flows one by one exposed a forwarding pitfall, idle flow meters, an NFStream export-ordering defect (median verdict delay 82.5\,s) and flow splitting at slow replay under unchanged macro-F1. \emph{Response:} blocking a host pair on its first attack verdict caused 60 false blocks in ten minutes; requiring 10 verdicts caused none in the ten live windows (3 in 5.9\,h with persistent counters) but missed the Mirai bots and the web attackers. Timed at flow export, verdicts let 11\% of attack flows start after their block, not 97\%. Code and results are public.
